% Options for packages loaded elsewhere
\PassOptionsToPackage{unicode}{hyperref}
\PassOptionsToPackage{hyphens}{url}
\documentclass[
]{article}
\usepackage{xcolor}
\usepackage[margin=1in]{geometry}
\usepackage{amsmath,amssymb}
\setcounter{secnumdepth}{-\maxdimen} % remove section numbering
\usepackage{iftex}
\ifPDFTeX
  \usepackage[T1]{fontenc}
  \usepackage[utf8]{inputenc}
  \usepackage{textcomp} % provide euro and other symbols
\else % if luatex or xetex
  \usepackage{unicode-math} % this also loads fontspec
  \defaultfontfeatures{Scale=MatchLowercase}
  \defaultfontfeatures[\rmfamily]{Ligatures=TeX,Scale=1}
\fi
\usepackage{lmodern}
\ifPDFTeX\else
  % xetex/luatex font selection
\fi
% Use upquote if available, for straight quotes in verbatim environments
\IfFileExists{upquote.sty}{\usepackage{upquote}}{}
\IfFileExists{microtype.sty}{% use microtype if available
  \usepackage[]{microtype}
  \UseMicrotypeSet[protrusion]{basicmath} % disable protrusion for tt fonts
}{}
\makeatletter
\@ifundefined{KOMAClassName}{% if non-KOMA class
  \IfFileExists{parskip.sty}{%
    \usepackage{parskip}
  }{% else
    \setlength{\parindent}{0pt}
    \setlength{\parskip}{6pt plus 2pt minus 1pt}}
}{% if KOMA class
  \KOMAoptions{parskip=half}}
\makeatother
\usepackage{graphicx}
\makeatletter
\newsavebox\pandoc@box
\newcommand*\pandocbounded[1]{% scales image to fit in text height/width
  \sbox\pandoc@box{#1}%
  \Gscale@div\@tempa{\textheight}{\dimexpr\ht\pandoc@box+\dp\pandoc@box\relax}%
  \Gscale@div\@tempb{\linewidth}{\wd\pandoc@box}%
  \ifdim\@tempb\p@<\@tempa\p@\let\@tempa\@tempb\fi% select the smaller of both
  \ifdim\@tempa\p@<\p@\scalebox{\@tempa}{\usebox\pandoc@box}%
  \else\usebox{\pandoc@box}%
  \fi%
}
% Set default figure placement to htbp
\def\fps@figure{htbp}
\makeatother
\setlength{\emergencystretch}{3em} % prevent overfull lines
\providecommand{\tightlist}{%
  \setlength{\itemsep}{0pt}\setlength{\parskip}{0pt}}
\usepackage{booktabs}
\usepackage{longtable}
\usepackage{array}
\usepackage{multirow}
\usepackage{wrapfig}
\usepackage{float}
\usepackage{colortbl}
\usepackage{pdflscape}
\usepackage{tabu}
\usepackage{threeparttable}
\usepackage{threeparttablex}
\usepackage[normalem]{ulem}
\usepackage{makecell}
\usepackage{xcolor}
\usepackage{bookmark}
\IfFileExists{xurl.sty}{\usepackage{xurl}}{} % add URL line breaks if available
\urlstyle{same}
\hypersetup{
  pdftitle={Assignment},
  hidelinks,
  pdfcreator={LaTeX via pandoc}}

\title{Assignment}
\author{}
\date{\vspace{-2.5em}}

\begin{document}
\maketitle

{
\setcounter{tocdepth}{4}
\tableofcontents
}
\subsection{Highway}\label{highway}

The traffic assignment process for the MRM is a multi-class ``user
equilibrium'' assignment, which is a multipath procedure where vehicle
trips are loaded from origin to destination through an iterative
process. During each iteration, the trips for each origin-destination
TAZ pair are assigned to a single shortest path along the network. The
loadings from the iterations are weighted in a manner that results, at
convergence, in the travel times along all paths being equal. This
ensures that no driver could improve the travel time by changing the
path. The MRM allows up to 250 iterations to ensure each assignment
reaches the set convergence of 0.0001. As part of the assignment, the
trip tables are assigned to the network.

\subsubsection{Highway Assignment Trip
Tables}\label{highway-assignment-trip-tables}

The assignment process is applied separately to each of the four time
periods (AM, PM, MI, NT). Then, those results are added together to
generate the total daily volume on each link.

For each of the time periods, six trip assignment tables are assigned:

\begin{enumerate}
\def\labelenumi{\arabic{enumi}.}
\item
  \textbf{Single-occupancy vehicles (SOV):} These trips do not have
  access to HOT lanes.
\item
  \textbf{Two-person carpool vehicle trips (HOV-2):} Vehicles with two
  persons have free access to HOT-2 express lanes (functional class 22
  and 82 for access). However, no HOT-2 lanes exist in the official
  model, as all existing and planned express lanes are HOT-3+.
\item
  \textbf{Three or more-person carpool vehicle trips (HOV-3):} Vehicles
  with three or more persons have free access to HOT-2 and HOT-3+ lanes
  (functional class 25 and 83 for access). HOT-3+ facilities in the
  model include the existing I-77 North Express Lanes, as well as the
  planned express lanes on I-77 South, US-74, and I-485 in build in
  future horizon years.
\item
  \textbf{Commercial vehicles:} These trips do not have access to HOT
  lanes during assignment.
\item
  \textbf{Medium trucks:} These trucks are assumed to have a
  passenger-car-equivalent (PCE) of 1.5 passenger cars when calculating
  congestion. These trips do not have access to HOT lanes or Peak Period
  Shoulder Lanes (PPSL's, functional class 24).
\item
  \textbf{Heavy trucks:} These trucks are assumed to have a PCE of 2.5
  when calculating congestion. These trips do not have access to HOT
  lanes or PPSL's.
\end{enumerate}

Trip tables from three sources are brought together for the assignment.
The mode choice program designates single-occupancy vehicle (SOV) and
carpool (HOV-2, HOV-3+) person trips for home-based trips (HBW, HBS,
HBO, SCH, HBU, JTW, ATW, NWK). Commercial vehicle, medium and heavy
trucks, and external trips (EI, IE) are brought in from trip
distribution. Through (EE) trips are brought in from the trip generation
step.

\subsubsection{Internal Trips}\label{internal-trips}

Carpool trips are divided into 2 person carpools and 3 or more person
carpools during mode choice. The number of 3+ carpools must be
``assigned'' an average number of persons per car. An average 3+ carpool
size is calculated for each trip purpose based on the 2023 household
travel survey .

\begingroup\fontsize{14}{16}\selectfont

\begin{longtabu} to \linewidth {>{\raggedright}X>{\raggedright}X}
\toprule
Purpose & Number of Persons\\
\midrule
HBW & 3.560 persons\\
HBO & 3.560 persons\\
HBU & 3.360 persons\\
HBSch & 3.560 persons\\
HBS & 3.560 persons\\
\addlinespace
ATW & 3.660 persons\\
\bottomrule
\end{longtabu}
\endgroup{}

\subsubsection{HOT Assignment}\label{hot-assignment}

The HOT assignment module is run after the normal feedback assignment
process is completed. The HOT lanes modeling process is applied by
time-of-day using the origin-destination matrices from the normal
feedback assignment process as the base set of trips. It runs the
assignment process five times for each time period (AM, PM, Midday, and
Night). Travel time from each round feeds back to a skimming process
that minimizes impedance for paths with or without using the HOT lanes.
The new travel time is averaged with the travel time from the previous
round to dampen oscillations between rounds.

After the skims are prepared, the origin-destination matrices are
updated with a ``Travel Time Savings'' core that is used to identify the
savings in travel time that are realized from using the HOT lanes. These
savings are used to identify diversion based on the amount of the HOT
toll and people's willingness to pay. Following the HOT assignment, a
new total assignment file (TOTASSN\_HOT.bin) is created that contains
the estimated link volumes. These volumes should be used in place of
TOTASSN.bin for analysis and model results reporting.

\subsubsection{Transit Assignment}\label{transit-assignment}

Transit Assignment in the MRM takes the transit production-attraction
output from the mode choice model and assigns trips to particular
transit route lines (bus, rail, etc.). Transit assignment differs from
the highway assignment in that it must consider the access and egress
modes in addition to the path.

The MRM uses a transit assignment method created by Caliper called
``Pathfinder''. The Pathfinder method can compute specific walk paths
for access and egress from any origin address to any destination address
and can assess the drive component for park and ride access on the
highway network and will yield the single best combination of park and
ride lot and transit service to the destination. To determine the best
path, transit assignment travel time is monetized so that it can be
combined with transit fares to yield a generalized cost. Considerations
other than transit fare and walk time include drive access time (park
and ride), transit route transfers, transit wait times, and timed
transfers. This allows transit assignment to have sensitivity to transit
fare increases, route schedule changes, and changes to the highway
network that impact transit route performance.

The transit assignment model assigns eighteen production-attraction
matrices, which include:

\begin{enumerate}
\def\labelenumi{\arabic{enumi}.}
\item
  \textbf{Peak Premium Walk Trips:} This trip table is created by adding
  walk-to-premium from peak HBW (all income groups), HBO (all income
  groups), HBU, and NHB models.
\item
  \textbf{Peak Premium Drive Trips:} This trip table is created by
  adding drive-to-premium trips from peak HBW (all income groups), HBO
  (all income groups), HBU, and NHB models.
\item
  \textbf{Peak Premium Drop-Off Trips:} This trip table is created by
  adding drop-to-premium trips from peak HBW (all income groups), HBO,
  HBU, and NHB modelsPeak Bus Walk Trips: This trip table is created by
  adding walk-to-bus trips from peak HBW (all income groups), HBO (all
  income groups), HBU, and NHB models.
\item
  \textbf{Peak Bus Drive Trips:} This trip table is created by adding
  drive-to-bus trips from peak HBW (all income groups), HBO (all income
  groups), HBU, and NHB models.
\item
  \textbf{Peak Bus Drop-Off Trips:} This trip table is created by adding
  drop-to-bus trips from peak HBW (all income groups), HBO (all income
  groups), HBU, and NHB models.
\item
  \textbf{Off-Peak Premium Walk Trips:} This trip table is created by
  adding walk-to-premium trips from off-peak HBW (all income groups),
  HBO (all income groups), HBU, and NHB models.
\item
  \textbf{Off-Peak Premium Drive Trips:} This trip table is created by
  adding drive-to-premium trips from off-peak HBW (all income groups),
  HBO (all income groups), HBU, and NHB models.
\item
  \textbf{Off-Peak Premium Drop-Off Trips:} This trip table is created
  by adding drop-to-premium trips from off-peak HBW (all income groups),
  HBO (all income groups), HBU, and NHB models.
\item
  \textbf{Off-Peak Bus Walk Trips:} This trip table is created by adding
  walk-to-bus trips from off-peak HBW (all income groups), HBO (all
  income groups), HBU, and NHB models.
\item
  \textbf{Off-Peak Bus Drive Trips:} This trip table is created by
  adding drive-to-bus trips from off-peak HBW (all income groups), HBO
  (all income groups), HBU, and NHB models.
\item
  \textbf{Off-Peak Bus Drop-Off Trips:} This trip table is created by
  adding drop-to-bus trips from off-peak HBW (all income groups), HBO
  (all income groups), HBU, and NHB models.
\end{enumerate}

\subsubsection{External Trips}\label{external-trips}

External trips consist of external-internal (EI) trips,
internal-external (IE,) trips, and external-external (EE) trips. These
three trip types comprise the external station counts of each facility
crossing the modeling area boundary. EI trips begin outside of the
modeling area and end inside of the modeling area at a TAZ. IE trips
begin inside the modeling area at TAZ and end outside the modeling area.
Both are captured at the external station as they exit or enter the
area. Finally, EE trips, commonly referred to as through trips, both
begin and end outside of the modeling area. These trips pass through the
modeling area entering at one external station and exiting at another.

Using vehicle classification counts we divided the vehicles crossing the
modeling area boundary at the external stations between autos,
commercial vehicles, medium trucks and heavy trucks. The home survey and
external roadside survey provided enough information to further divide
IE and EI auto trips, respectively, into work and non-work trips from
productions to attractions. However, we were not able to obtain enough
information from the external roadside survey to do the same for EI
trips from attractions to productions. As such, we applied the IE
attraction to production proportion obtained from the home survey to the
EI attraction to productions trips.

In trip distribution, IE and EI work and non-work trips are modeled as
person trips rather than vehicle trips to allow for the possible
inclusion in mode choice at some future date.

\subsubsection{Volume Delay Functions}\label{volume-delay-functions}

The MRM uses the commonly used Bureau of Public Roads (BPR) formula as
the volume-delay function to relate link travel time to the
volume/capacity ratio. The BPR formula is shown below.

The basic BPR function is expressed as:

\[
t = t_0 \times \left( 1 + \alpha \times \left( \frac{v}{c} \right)^\beta \right)
\]

Where:

\begin{itemize}
\tightlist
\item
  \textbf{t}: Travel time on the link\\
\item
  \textbf{t₀}: Free flow travel time\\
\item
  \textbf{v}: Traffic volume on the link\\
\item
  \textbf{c}: Capacity of the link\\
\item
  \textbf{α} and \textbf{β}: Calibration parameters, usually determined
  based on specific road characteristics
\end{itemize}

Different functionally classified roads are known to have different
alpha and beta coefficients which describe the interaction between
travel behavior, congestion, and speed on a particular corridor. \#\#\#
Function Calibration Parameters

\begingroup\fontsize{14}{16}\selectfont

\begin{longtabu} to \linewidth {>{\raggedright\arraybackslash}p{10em}>{\raggedright}X>{\raggedright}X>{\raggedright}X>{\raggedright}X}
\toprule
Functional Classification & Area Type & Lanes & Alpha & Beta\\
\midrule
\addlinespace[0.3em]
\multicolumn{5}{l}{\textbf{Freeway}}\\
\hspace{1em}1 & 2-3 &  & 0.83 & 10.00\\
\hspace{1em}1 & 1 &  & 0.7470 & 9.00\\
\hspace{1em}1 & 4 &  & 0.85 & 10.00\\
\hspace{1em}1 & 5 &  & 0.88 & 10.30\\
\addlinespace[0.3em]
\multicolumn{5}{l}{\textbf{Expressway}}\\
\hspace{1em}\hspace{1em}2 & 1-3 &  & 0.83 & 8.00\\
\hspace{1em}2 & 4 &  & 0.83 & 8.00\\
\hspace{1em}2 & 5 &  & 0.8549 & 8.24\\
\addlinespace[0.3em]
\multicolumn{5}{l}{\textbf{Class II}}\\
\hspace{1em}\hspace{1em}3 & 1-4 &  & 0.62 & 7.13\\
\hspace{1em}3 & 5 &  & 0.6386 & 7.3439\\
\hspace{1em}3 & 1-3 & 3 & 0.70 & 7.50\\
\hspace{1em}3 & 4 & 3 & 0.65 & 7.50\\
\hspace{1em}3 & 5 & 3 & 0.67 & 7.73\\
\addlinespace[0.3em]
\multicolumn{5}{l}{\textbf{Major Arterial}}\\
\hspace{1em}4 & 1-2 &  & 1.00 & 5.00\\
\hspace{1em}4 & 3-4 &  & 0.60 & 5.00\\
\hspace{1em}4 & 5 &  & 0.68 & 6.85\\
\hspace{1em}4 & 3-4 & 3 & 0.65 & 5.50\\
\hspace{1em}4 & 5 & 3 & 0.72 & 7.21\\
\addlinespace[0.3em]
\multicolumn{5}{l}{\textbf{Minor Arterial}}\\
\hspace{1em}\hspace{1em}5 & 1-2 &  & 0.75 & 3.00\\
\hspace{1em}5 & 3-4 &  & 0.65 & 4.00\\
\hspace{1em}5 & 5 &  & 0.73 & 4.89\\
\hspace{1em}5 & 1-2 & 3 & 0.75 & 2.10\\
\hspace{1em}5 & 3-4 & 3 & 0.65 & 4.00\\
\hspace{1em}5 & 5 & 3 & 0.77 & 4.85\\
\addlinespace[0.3em]
\multicolumn{5}{l}{\textbf{Collector}}\\
\hspace{1em}6 & 1-4 &  & 0.76 & 3.80\\
\hspace{1em}6 & 5 &  & 0.73 & 3.91\\
\hspace{1em}6 & 1-4 & 3 & 0.80 & 4.00\\
\hspace{1em}6 & 5 & 3 & 0.82 & 4.12\\
\addlinespace[0.3em]
\multicolumn{5}{l}{\textbf{Local}}\\
\hspace{1em}7 & 1-4 &  & 0.95 & 3.80\\
\hspace{1em}7 & 5 &  & 0.98 & 3.91\\
\addlinespace[0.3em]
\multicolumn{5}{l}{\textbf{Ramp}}\\
\hspace{1em}8 & 1-4 &  & 0.79 & 6.65\\
\hspace{1em}8 & 5 &  & 0.8137 & 6.8495\\
8 & 5 & 3 & 0.81 & 6.85\\
\addlinespace[0.3em]
\multicolumn{5}{l}{\textbf{Freeway Ramp}}\\
\hspace{1em}9 & 1-4 &  & 0.79 & 9.50\\
\hspace{1em}9 & 5 &  & 0.81 & 9.79\\
\addlinespace[0.3em]
\multicolumn{5}{l}{\textbf{HOT2- Freeway / Toll Only Lanes}}\\
\hspace{1em}22/23 & 1-3 &  & 0.83 & 10.00\\
\hspace{1em}22/23 & 4 &  & 0.85 & 10.00\\
\hspace{1em}22/23 & 5 &  & 0.88 & 10.30\\
\addlinespace[0.3em]
\multicolumn{5}{l}{\textbf{Peak Period Shoulder Lanes}}\\
\hspace{1em}\hspace{1em}24 & 1 &  & 0.83 & 10.00\\
\hspace{1em}24 & 2-3 &  & 0.62 & 7.13\\
\hspace{1em}24 & 4 &  & 0.85 & 10.00\\
\hspace{1em}24 & 5 &  & 0.88 & 10.30\\
\hspace{1em}24 & 2-3 & 3 & 0.70 & 7.50\\
\addlinespace[0.3em]
\multicolumn{5}{l}{\textbf{HOT3- Freeway}}\\
\hspace{1em}25 & 1 &  & 0.83 & 10.00\\
\hspace{1em}25 & 2-3 &  & 0.60 & 5.00\\
\hspace{1em}25 & 4 &  & 0.85 & 10.00\\
\hspace{1em}25 & 5 &  & 0.88 & 10.30\\
\hspace{1em}25 & 2-3 & 3 & 0.65 & 5.50\\
\addlinespace[0.3em]
\multicolumn{5}{l}{\textbf{HOT2 \& HOT3 Access}}\\
\hspace{1em}\hspace{1em}82 &  &  & 0.79 & 9.50\\
83 &  &  & 0.79 & 9.50\\
\bottomrule
\end{longtabu}
\endgroup{}

\end{document}
